% Imporditud: tools/import_eksperimendid.py
% Allikas: Eesti füüsikaolümpiaadi eksperimendi ülesannete
%   kogu (Rael Kalda, 2023), ülesanne Ü22, lahendus L22.
\setAuthor{Erkki Tempel}
\setRound{piirkonnavoor}
\setYear{2015}
\setNumber{P E1}
\setDifficulty{1}
\setTopic{vedelike mehaanika}
\setVahendid{metallist keha, dünamomeeter, anum veega ($\rho$v = 1 g/cm$^3$)}
\setPunktid{10}
\setAllikas{Eksperimendi kogu Ü22, lahendus lk 41}
\setLahendusseis{toores}

\prob{Keha ruumala}
Määrake metallist keha ruumala.

\hint

\solu
Riputame keha dünamomeetri külge, saame teada kehale mõjuva raskusjõu $N_1$ = Fr, kus $N_1$ on dünamomeetri näit. Riputame keha dünamomeetri külge ning sukeldades vette, näitab dünamomeeter ($N_2$) keha raskusjou Fr ja üleslükkejõu Fü = $\rho$vgV vahet

$N_2$ = Fr $-$ Fü = $N_1$ $- \rho$vgV $\Rightarrow$ V = $N_1$ $-$ $N_2$ . $\rho$v g

Hindamine: Eksperimendi idee --- 3 p. Raskusjõu mõõtmine --- 1 p. Vette asetatud kehale mõjuva jõu mõõtmine --- 1 p. Seos $N_2$ = Fr $-$ Fü = $N_1$ $- \rho$vgV annab 1 p. Keha ruumala V avaldamine --- 2 p. Kordusmõõtmised (3 katset) --- 1 p. Vastus erineb õigest tulemusest kuni 5 cm$^3$ --- 1 p..
\probend
